%% Macros de dados do documento
\autor{Henrique Machado de Oliveira}
\titulo{Redes Neurais Convolucionais para o Controle de Qualidade de Rochas Ornamentais: Da Classificação à Segmentação de Anomalias utilizando pipeline com abordagem Teacher-Student}

%%%Sistema Hierárquico de Inteligência Artificial para o Controle de Qualidade de Rochas Ornamentais: Da Classificação à Segmentação de Irregularidades
%%%Análise Comparativa entre Modelos de Segmentação Generalistas e Arquiteturas em Cascata para Identificação de Defeitos em Rochas Ornamentais
%%%Otimização da Segmentação de Superfícies Pétreas através de Classificação Prévia e Redes Especialistas
%%%Abordagem Hierárquica em Deep Learning: Impacto da Especialização de Modelos na Precisão da Segmentação de Rochas Ornamentais

\instituicao{Instituto Federal do Espírito Santo}
\curso{Bacharelado em Sistemas de Informação}
\local{Cachoeiro de Itapemirim}
\data{2026}
\tipotrabalho{Trabalho de Graduação}
\preambulo{Trabalho de Conclusão de Curso apresentado à Coordenadoria do Curso de Sistemas de Informação do Instituto Federal de Educação, Ciência e Tecnologia do Espírito Santo, como requisito parcial para a obtenção do título de Bacharel em Sistemas de Informação.}

\orientador[Orientador][Prof. Dr.]{Rafael}[Silva Guimarães]{Instituto Federal do Espírito Santo}
%%%\coorientador[Coorientador][Prof. Dr.]{Everson}[Scherrer Borges]{Instituto Federal do Espírito Santo}
%%%\coorientador[Coorientador][Prof. Dr.]{Igor}[Henrique Beloti Pizetta]{Instituto Federal do Espírito Santo}

\examinadori{Profa. Dra. Fulana de Tal}{Instituto Federal do Espírito Santo}{Examinadora}
\examinadorii{Prof. Dr. Cicrano de Tal}{Instituto Federal do Espírito Santo}{Examinador}

% TODO: data real da defesa, a confirmar com o orientador (ver docs/pendencias.md).
\approvaldate{15}{Outubro}{2026}

% Palavras-chaves
\palavraschave{Visão Computacional}[Segmentação Semântica][Rochas Ornamentais][Modelos de Fundação][Pseudo-rotulagem]
\keywords{Computer Vision}[Semantic Segmentation][Dimension Stone][Foundation Models][Pseudo-labeling]

% Ficha catalográfica
\fichabiblioteca{
    Dados Internacionais de Catalogação-na-Publicação (CIP) \\
    (Biblioteca Nilo Peçanha do Instituto Federal do Espírito Santo)
}
\fichaautor{Elaborada por XXXXXXXXXXXXXXXXXXXXX – CRB-X/ES - XXX}
\fichatags{
    1. Visão computacional.  
    2. Aprendizado profundo.  
    3. Segmentação de imagens.  
    4. Rochas ornamentais – Controle de qualidade.  
    5. Inteligência artificial.  
    I. Guimarães, Rafael Silva.  
    II. Instituto Federal do Espírito Santo.  
    III. Título.
}

\cutter{X999y}
\cdd{000.00}
\selectlanguage{brazil}

% Arial (para usar times-new-roman, comente as três linhas abaixo)
\usepackage{helvet}
\renewcommand{\familydefault}{\sfdefault}
\renewcommand{\ABNTEXchapterfont}{\bfseries}
% Comandos para revisar o texto
\usepackage{easyReview} 


%% Elementos opcionais
\editardedicatoria{A dedicatória é um elemento opcional. Contém o oferecimento do trabalho adeterminada pessoa ou a pessoas (ASSOCIAÇÃO BRASILEIRA DE NORMAS TÉCNICAS, 2011b).}
\editaragradecimentos{A um elemento opcional. Localiza-se após a folha de aprovação e deve ser dirigido àqueles que realmente contribuíram, de maneira relevante, para a elaboração do trabalho. Deve-se utilizar uma linguagem simples (ASSOCIAÇÃO BRASILEIRA DE NORMAS TÉCNICAS, 2011b).}
\editarepigrafe{\textbf{Epigrafe:} É um elemento opcional. É uma citação relacionada ao assunto do trabalho desenvolvido, seguida da indicação de autoria (ASSOCIAÇÃO BRASILEIRA DE NORMAS TÉCNICAS, 2011b).   Deve-se seguir as regras do uso da citação NBR 10.520/2002.}


%% Listas [Siglas e Simbolos]
\editarlistasiglas{\begin{siglas}
    \item[ARIA] Análise e Reconhecimento Inteligente de Anomalias
    \item[CLIP] \textit{Contrastive Language--Image Pre-training}
    \item[CNN] Rede Neural Convolucional (\textit{Convolutional Neural Network})
    \item[GPU] Unidade de Processamento Gráfico (\textit{Graphics Processing Unit})
    \item[Ifes] Instituto Federal do Espírito Santo
    \item[IoU] \textit{Intersection over Union}
    \item[mAP] \textit{mean Average Precision}
    \item[NMS] \textit{Non-Maximum Suppression}
    \item[SAM] \textit{Segment Anything Model}
    \item[YOLO] \textit{You Only Look Once}
\end{siglas}
}
\editarlistasimbolos{\begin{simbolos}
    \item[$\lambda$] Letra grega labda
\end{simbolos}}